% ============================================================
% Section 11: Early 20th Century (Slides 79-83)
% ============================================================
\section{Early 20th Century}

% ----------------------------------------------------------
% Slide 79: Section Divider
% ----------------------------------------------------------
{
\setbeamercolor{background canvas}{bg=digitalBlue}
\begin{frame}[plain]
  \centering
  \vspace{2cm}
  \textcolor{white}{\Huge\bfseries Crisis, War,}\\[0.4em]
  \textcolor{white}{\Huge\bfseries and Computers}\\[1em]
  \textcolor{white!80}{\large 1900--1950 CE}\\[1.5em]
  \textcolor{white!60}{\normalsize Hilbert \(\cdot\) G\"odel \(\cdot\) Turing \(\cdot\) von Neumann \(\cdot\) Nash}

  \note{
    Transition: ``The 20th century opened with Hilbert's challenge: can we prove
    math is consistent? The answer was shocking.''
    This section covers the most turbulent half-century in mathematics ---
    foundations crises, world wars, and the birth of the computer.
  }
\end{frame}
}

% ----------------------------------------------------------
% Slide 80: David Hilbert and the Foundations Crisis
% ----------------------------------------------------------
\begin{frame}{David Hilbert --- The Foundations Crisis}

  \begin{columns}[T]
    \begin{column}{0.5\textwidth}
      \begin{personbox}[David Hilbert (1862--1943)]{digitalBlue}
        \textbf{Origin:} K\"onigsberg, Prussia (modern Kaliningrad, Russia);
        worked in G\"ottingen\\[3pt]
        \textbf{Key work:} In 1900, posed \alert{23 unsolved problems} that
        shaped 20th-century mathematics. Hilbert's program aimed to put
        \emph{all} of mathematics on a provably consistent, complete
        foundation.
      \end{personbox}

      \vspace{0.4em}
      \begin{personbox}[Kurt G\"odel (1906--1978)]{digitalBlue}
        \textbf{Origin:} Brno, Austria-Hungary (now Czechia); worked at
        Princeton, USA\\[3pt]
        \textbf{Incompleteness Theorems (1931):} Any sufficiently powerful
        mathematical system contains true statements that
        \textbf{cannot be proved} within the system.
      \end{personbox}
    \end{column}

    \begin{column}{0.47\textwidth}
      % Hilbert's 23 problems (selected)
      \visualplaceholder[5cm]{G "odel's theorem in plain English}
    \end{column}
  \end{columns}

  \note{
    ``Hilbert asked: can we prove that mathematics is free of contradictions?
    G\"odel answered: no. And he proved it mathematically.''
    This was a philosophical earthquake. Mathematics --- the most certain
    of all sciences --- had fundamental limits.
    G\"odel's proof is one of the most profound results in all of intellectual history.
  }
\end{frame}

% ----------------------------------------------------------
% Slide 81: Alan Turing -- The Birth of Computer Science
% ----------------------------------------------------------
\begin{frame}{Alan Turing --- The Birth of Computer Science}

  \begin{columns}[T]
    \begin{column}{0.48\textwidth}
      \begin{personbox}[Alan Mathison Turing (1912--1954)]{digitalBlue}
        \textbf{Origin:} London, England; worked at Cambridge, Bletchley Park,
        Manchester\\[3pt]
        \textbf{Key work:}
        \begin{itemize}\setlength\itemsep{1pt}
          \item The \alert{Turing machine} (1936) --- theoretical foundation
                of all computers
          \item Proved the \textbf{halting problem} is unsolvable
          \item Broke wartime Enigma at Bletchley Park, building on
                \textbf{Marian Rejewski} (1905--1980, Poland) and
                Polish cryptanalysts who cracked earlier Enigma versions
                in the 1930s
          \item Pioneered the ``Turing test'' for AI
        \end{itemize}
      \end{personbox}
    \end{column}

    \begin{column}{0.49\textwidth}
      % TikZ: Turing machine diagram
      \visualplaceholder[5cm]{Tape}
    \end{column}
  \end{columns}

  \note{
    ``Polish mathematicians cracked the first version of Enigma. Turing and his
    team broke the far more complex wartime versions. Every computer, every phone,
    every AI runs on principles Turing laid down in a 1936 paper.''
    \verifytag{Polish contribution well-documented. Rejewski broke initial Enigma
    cipher in 1932; Poles shared work with UK/France in July 1939. Turing and team
    then broke more complex wartime versions.}
    Turing was persecuted for being gay and died at 41.
    He was posthumously pardoned in 2013.
  }
\end{frame}

% ----------------------------------------------------------
% Slide 82: John von Neumann
% ----------------------------------------------------------
\begin{frame}{John von Neumann --- The Last Great Polymath}

  \begin{columns}[T]
    \begin{column}{0.47\textwidth}
      \begin{personbox}[John von Neumann (1903--1957)]{digitalBlue}
        \textbf{Born:} Neumann J\'anos Lajos\\
        \textbf{Origin:} Budapest, Hungary; worked in Germany and
        the USA (Princeton)\\[3pt]
        \textbf{Key work:} Contributions to quantum mechanics, game theory,
        computer architecture, cellular automata, nuclear weapons design,
        and set theory axiomatisation. Possibly the last great polymath.
      \end{personbox}

      \vspace{0.3em}
      \begin{insightbox}
        The computer you use, the phone in your pocket --- they all follow
        the architecture he designed in the 1940s.
      \end{insightbox}
    \end{column}

    \begin{column}{0.5\textwidth}
      % TikZ: Von Neumann architecture
      \visualplaceholder[5cm]{CPU block}
    \end{column}
  \end{columns}

  \note{
    ``The computer you use, the phone in your pocket --- they all follow the
    architecture he designed in the 1940s.''
    von Neumann had an almost superhuman memory and processing speed.
    He contributed to so many fields that it is hard to summarise his work
    in one slide.
    The key insight of his architecture: store both program and data
    in the same memory (the ``stored-program concept'').
  }
\end{frame}

% ----------------------------------------------------------
% Slide 83: John Nash -- Game Theory
% ----------------------------------------------------------
\begin{frame}{John Nash --- Game Theory and the Nash Equilibrium}

  \begin{columns}[T]
    \begin{column}{0.48\textwidth}
      \begin{personbox}[John Forbes Nash Jr.\ (1928--2015)]{digitalBlue}
        \textbf{Origin:} Bluefield, West Virginia, USA; worked at MIT
        and Princeton\\[3pt]
        \textbf{Key work:}
        \begin{itemize}\setlength\itemsep{2pt}
          \item \alert{Nash equilibrium} (1950) --- no player can improve
                their outcome by unilaterally changing strategy
          \item Applications: economics, evolutionary biology, AI
          \item Nobel Prize in Economics (1994)
          \item Struggled with schizophrenia for decades; eventually recovered
        \end{itemize}
      \end{personbox}

      \vspace{0.3em}
      {\small\itshape Portrayed in ``A Beautiful Mind'' (2001)}
    \end{column}

    \begin{column}{0.49\textwidth}
      % TikZ: Nash equilibrium payoff matrix
      \visualplaceholder[5cm]{Matrix frame}
    \end{column}
  \end{columns}

  \note{
    ``Every time countries negotiate treaties, companies set prices, or AI agents
    make decisions --- game theory is at work. And its foundations were laid by a
    mathematician who nearly lost everything to mental illness.''
    Nash's work is used today in everything from auction design to machine learning
    multi-agent systems.
    His recovery from schizophrenia after decades of illness is extraordinary.
    He and his wife Alicia died in a taxi accident in 2015.
  }
\end{frame}
